\subsection{Point distributions and tangent spaces}

In this section, X denotes a manifold of class C \(^{r}\) .

13.3.1. Let \(x \in X\) and let \(k\) be an integer \(\leqslant r\) . Let \(c = (U, \varphi, E)\) a chart of \(X\) centered at \(x\) ; the isomorphism \(\theta_c \colon \mathsf{TS}^{(r)}(E) \to \mathsf{T}_x^{(r)}(X)\) defined in 13.2.1 maps \(\mathsf{TS}^{(k)}(E)\) onto \(\mathsf{T}_x^{(k)}(X)\) and \(\mathsf{TS}^{(k-1)}(E)\) onto \(\mathsf{T}_x^{(k-1)}(X)\) ; by restriction and passage to the quotient, it induces an isomorphism

\[
\theta_ {c, k} \colon \mathsf {T S} ^ {k} (\mathrm{E}) = \mathsf {T S} ^ {(k)} (\mathrm{E}) / \mathsf {T S} ^ {(k - 1)} (\mathrm{E}) \rightarrow \mathsf {T} _ {x} ^ {(k)} (\mathrm{X}) / \mathsf {T} _ {x} ^ {(k - 1)} (\mathrm{X}).
\]

On the other hand, c defines an isomorphism, already denoted \(\theta_{c}\) , from E onto the tangent space \(T_{x}(X)\) , cf. 5.5.1. Let \(i_{k}\) denote the composite

\[
\mathrm{T} _ {x} ^ {(k)} (\mathrm{X}) / \mathrm{T} _ {x} ^ {(k - 1)} (\mathrm{X}) \xrightarrow {\theta_ {c , k} ^ {- 1}} \mathrm{TS} ^ {k} (\mathrm{E}) \xrightarrow {\mathrm{TS} ^ {k} (\theta_ {c})} \mathrm{TS} ^ {k} (\mathrm{T} _ {x} (\mathrm{X})).
\]

The isomorphism \(i_k\) is independent of the choice of \(c\) ; for \(k = 1\) , it is used to identify \(\mathrm{T}_x^{(1)+}(\mathrm{X}) = \mathrm{T}_x^{(1)}(\mathrm{X}) / \mathrm{T}_x^{(0)}(\mathrm{X})\) with the tangent space \(\mathrm{T}_x(\mathrm{X})\) . If \(t \in \mathrm{T}_x^{(k)}(\mathrm{X})\) , one may write \(i_k(t)\) for the image under \(i_k\) of the class of \(t\) modulo \(\mathrm{T}_x^{(k-1)}(\mathrm{X})\) .

Set

\[
\operatorname{gr} _ {k} \mathrm{T} _ {x} ^ {(r)} (\mathrm{X}) = \mathrm{T} _ {x} ^ {(k)} (\mathrm{X}) / \mathrm{T} _ {x} ^ {(k - 1)} (\mathrm{X}) \quad \text { and } \quad \operatorname{gr} \mathrm{T} _ {x} ^ {(r)} (\mathrm{X}) = \bigoplus_ {0 \leqslant k \leqslant r} \operatorname{gr} _ {k} \mathrm{T} ^ {(r)} (\mathrm{X}).
\]

One says that gr \(\mathbf{T}_{x}^{(r)}(\mathbf{X})\) is the graded object associated with the increasing filtration \((\mathbf{T}_{x}^{(k)}(\mathbf{X}))_{k\leq r}\) of \(\mathbf{T}_{x}^{(r)}(\mathbf{X})\) . The \(i_k\) define an isomorphism of graded vector spaces

\[
i \colon \operatorname{gr} \mathrm{T} _ {x} ^ {(r)} (\mathrm{X}) \rightarrow \mathrm{TS} ^ {(r)} (\mathrm{T} _ {x} (\mathrm{X})).
\]

If \(r = \infty\) or \(\omega\) , \(i\) is an isomorphism of gr \(T_x^{(r)}(X)\) onto \(TS(T_x(X))\) . When one wishes to specify \(x\) (or X, or both), one writes \(i_x\) (or \(i_X\) , or \(i_{X,x}\) ) in place of \(i\) .

13.3.2. In addition to the above hypotheses, suppose that X is locally of finite dimension. The isomorphisms \(i_{k}\) relative to the various points of X define an isomorphism of vector bundles

\[
i _ {k} \colon \mathrm{T} ^ {(k)} (\mathrm{X}) / \mathrm{T} ^ {(k - 1)} (\mathrm{X}) \rightarrow \mathrm{TS} ^ {k} (\mathrm{T} (\mathrm{X}))
\]

where TS \(^{k}\) (T(X)) denotes the vector bundle of class C \(^{r-1}\) deduced from T(X) by the finite-dimensional vector functor TS \(^{k}\) (cf. 7.6.5). For \(k \geqslant 1\) , \(i_{k}\) is an isomorphism of class C \(^{r-k}\) ; for \(k = 0\) , it is the identity map of the trivial bundle \(\mathrm{K}_{\mathrm{X}}\) .

13.3.3. Example. --- With hypotheses as above, let \(\xi = (\xi^{1}, \ldots, \xi^{n})\) be a coordinate system of X at \(x\) . Let \((\partial_{i,x})_{1 \leqslant i \leqslant n}\) denote the basis of \(\mathrm{T}_{x}(\mathrm{X})\) defined by \(\xi\) (cf. 5.5.8), and let \(((\Delta_{\xi}^{\infty})_{x})_{|\alpha| \leqslant k}\) denote the basis of \(\mathrm{T}_{x}^{(k)}(\mathrm{X})\) defined in 13.2.6. One has:

\[
\begin{array}{l l} i _ {k} ((\Delta_ {\xi} ^ {\alpha}) _ {x}) = 0 & \text {if} \quad | \alpha | <   k \\ i _ {k} ((\Delta_ {\xi} ^ {\alpha}) _ {x}) = \gamma_ {\alpha_ {1}} (\partial_ {1, x}) \dots \gamma_ {\alpha_ {n}} (\partial_ {n, x}) & \text {if} \quad | \alpha | = k. ^ {(1)} \end{array}
\]

\(^{1}\) Here again, the product of the \(\gamma_{\alpha_{i}}(\partial_{i,x})\) is a symmetric product in \(\mathrm{TS}(\mathrm{T}_{x}(\mathrm{X}))\).

13.3.4. Suppose that K has characteristic zero or that X is locally of finite dimension. Let k be an integer such that \(0 \leqslant k \leqslant r\) and let \(x \in X\) ; let F be a Banach space. By 12.6.8, one has an exact sequence

\[
0 \rightarrow \mathrm{P} _ {k} (\mathrm{T} _ {x} (\mathrm{X}); \mathrm{F}) \stackrel {{\iota}} {{\rightarrow}} \mathrm{P} _ {x} ^ {k} (\mathrm{F} _ {\mathrm{X}}) \stackrel {{\sigma}} {{\rightarrow}} \mathrm{P} _ {x} ^ {k - 1} (\mathrm{F} _ {\mathrm{X}}) \rightarrow 0,\tag{i}
\]

with \(\sigma = r^{k, k-1}\) . On the other hand, by 13.3.1, one has an exact sequence

\[
0 \leftarrow \mathsf {T S} ^ {k} (\mathrm{T} _ {x} (\mathrm{X})) \stackrel {{i}} {{\leftarrow}} \mathrm{T} _ {x} ^ {(k)} (\mathrm{X}) \stackrel {{s}} {{\leftarrow}} \mathrm{T} _ {x} ^ {(k - 1)} (\mathrm{X}) \leftarrow 0,\tag{ii}
\]

where s is the inclusion of \(\mathrm{T}_{x}^{(k-1)}(\mathrm{X})\) to \(\mathrm{T}_{x}^{(k)}(\mathrm{X})\) and \(i=i_{k}\) . These exact sequences are "paired in F". More precisely, write them in the form:

\begin{enumerate}
\def\labelenumi{(\roman{enumi})}
\tightlist
\item
\end{enumerate}

\[
0 \rightarrow \mathrm{A} \xrightarrow {\iota} \mathrm{B} \xrightarrow {\sigma} \mathrm{C} \rightarrow 0\tag{ii}
\]

\[
0 \leftarrow A ^ {\prime} \stackrel {i} {\leftarrow} B ^ {\prime} \stackrel {s} {\leftarrow} C ^ {\prime} \leftarrow 0.
\]

If \(a \in \mathbf{A}\) and \(a' \in \mathbf{A}'\) , the element \(\langle a, a' \rangle\) of F is defined by 13.1.2; if \(b \in \mathbf{B}\) and \(b' \in \mathbf{B}'\) , (resp. if \(c \in \mathbf{C}\) and \(c' \in \mathbf{C}'\) ), the element \(\langle b, b' \rangle\) (resp. \(\langle c, c' \rangle\) ) of F is defined by 13.2.2; one then has:

\[
\langle \iota a, b ^ {\prime} \rangle = \langle a, i b ^ {\prime} \rangle \quad \text { and } \quad \langle \sigma b, c ^ {\prime} \rangle = \langle b, s c ^ {\prime} \rangle \quad \text { if } \quad a \in \mathbf {A}, \quad b \in \mathbf {B}, \quad b ^ {\prime} \in \mathbf {B} ^ {\prime}, \quad c ^ {\prime} \in \mathbf {C} ^ {\prime}.
\]

13.3.5. Let Y be a manifold of class \(C^{r}\) let \(\varphi: X \to Y\) be a morphism of class \(C^{r}\) and let \(x \in X\) , \(y \in Y\) such that \(y = \varphi(x)\) . The map \(\varphi_{*}: T_{x}^{(r)}(X) \to T_{y}^{(r)}(Y)\) defined in 13.2.3 is compatible with the filtrations of these spaces; it defines, by passage to the associated graded objects, a linear map

\[
\operatorname{gr} \left(\varphi_ {*}\right): \operatorname{gr} \mathrm{T} _ {x} ^ {(r)} (\mathrm{X}) \rightarrow \operatorname{gr} \mathrm{T} _ {y} ^ {(r)} (\mathrm{Y}).
\]

Let us also denote by \(\mathsf{TS}^{(r)}(\mathsf{T}_x(\varphi))\) the map from \(\mathsf{TS}^{(r)}(\mathsf{T}_x(\mathbf{X}))\) to \(\mathsf{TS}^{(r)}(\mathsf{T}_y(\mathbf{Y}))\) induced by the canonical extension \(\mathsf{TS}(\mathsf{T}_x(\varphi))\) of \(\mathsf{T}_x(\varphi)\) . The diagram

\[
\begin{array}{c} \operatorname{gr} T _ {x} ^ {(r)} (X) \xrightarrow {\operatorname{gr} (\varphi_ {*})} \operatorname{gr} T _ {y} ^ {(r)} (Y) \\ \Biggl \downarrow i _ {\mathbf {X}, x} \\ \operatorname{TS} ^ {(r)} (T _ {x} (X)) \xrightarrow {\operatorname{TS} ^ {(r)} (T _ {x} (\varphi))} \operatorname{TS} ^ {(r)} (T _ {y} (Y)) \end{array}
\]

is commutative.

13.3.6. Suppose that K is of characteristic zero. Using the isomorphisms \(\varphi_{\mathbb{M}}: \mathsf{S}^k(\mathbf{M}) \to \mathsf{TS}^k(\mathbf{M})\) defined in A, IV, § 5, no. 8, one may replace, in all that precedes, the \(\mathsf{TS}^k(\mathbf{T}_x(\mathbf{X}))\) by the symmetric powers \(k\) th \(\mathsf{S}^k(\mathbf{T}_x(\mathbf{X}))\) of the tangent space \(\mathbf{T}_x(\mathbf{X})\) .

